\section{Evaluation}
\label{sec:evaluation}

\subsection{Headline metrics}

We evaluate all eight (model, scenario) combinations on the held-out \num{2303236}-row test set covering June~2024. The headline metric is precision-recall area under curve (PR-AUC); F1 at the tuned threshold is reported alongside. \Cref{tab:headline-metrics} lays out the eight numbers; \cref{fig:metrics-summary,fig:pr-roc} render them visually.

\begin{table}[htbp]
    \centering
    \caption{Test-set PR-AUC and F1 for the eight (model, scenario) combinations. Higher is better in both columns. The right-most column reports the absolute PR-AUC gain when the same model moves from scenario~A (pre-departure) to scenario~B (inflight). XGBoost wins both scenarios; the inflight uplift is dominated by a single feature (see \cref{sec:xai}).}
    \label{tab:headline-metrics}
    \small
    \begin{tabular}{l
        S[table-format=1.3] S[table-format=1.3]
        S[table-format=1.3] S[table-format=1.3]
        S[table-format=1.3]}
        \toprule
        & \multicolumn{2}{c}{Scenario A (pre-departure)} & \multicolumn{2}{c}{Scenario B (inflight)} & \\
        \cmidrule(lr){2-3}\cmidrule(lr){4-5}
        Model & {PR-AUC} & {F1} & {PR-AUC} & {F1} & {A$\to$B uplift} \\
        \midrule
        Logistic Regression  & 0.280 & 0.352 & 0.541 & 0.546 & 0.261 \\
        LightGBM             & 0.520 & 0.525 & 0.951 & 0.903 & 0.431 \\
        \textbf{XGBoost}     & {\bfseries 0.552} & {\bfseries 0.530} & {\bfseries 0.952} & {\bfseries 0.906} & 0.400 \\
        GraphSAGE            & 0.520 & 0.518 & 0.921 & 0.891 & 0.401 \\
        \bottomrule
    \end{tabular}
\end{table}

\begin{figure}[htbp]
    \centering
    \begin{subfigure}{0.48\textwidth}
        \centering
        \includegraphics[width=\linewidth]{stop_level/metrics_summary_A.png}
        \caption{Scenario A (pre-departure).}
    \end{subfigure}\hfill
    \begin{subfigure}{0.48\textwidth}
        \centering
        \includegraphics[width=\linewidth]{stop_level/metrics_summary_B.png}
        \caption{Scenario B (inflight).}
    \end{subfigure}
    \caption{Per-model PR-AUC, F1, precision and recall on the test set. The three tree- and graph-based models cluster tightly in scenario~A; in scenario~B they all reach PR-AUC above \num{0.92}, while logistic regression saturates around \num{0.54}.}
    \label{fig:metrics-summary}
\end{figure}

\begin{figure}[htbp]
    \centering
    \includegraphics[width=0.92\textwidth]{stop_level/pr_roc_grid.png}
    \caption{ROC curves (left column) and precision-recall curves (right column) for all four models on each scenario. In scenario~B the precision-recall curves hug the top-right corner: the A$\to$B PR-AUC jump from roughly \num{0.52} to \num{0.95} is visible at a glance.}
    \label{fig:pr-roc}
\end{figure}

\subsection{Calibration}

A high PR-AUC is necessary but not sufficient: a downstream operations team will only use the predicted probabilities directly if those probabilities mean what they say. \Cref{fig:calibration} shows reliability diagrams for each (model, scenario) combination. LightGBM and XGBoost stay close to the perfect-calibration diagonal in scenario~B; logistic regression is the worst on both scenarios, with an Expected Calibration Error (ten-bin quantile) of \num{0.319} versus \num{0.043} for XGBoost.

\begin{figure}[htbp]
    \centering
    \includegraphics[width=0.92\textwidth]{stop_level/calibration_grid.png}
    \caption{Reliability diagrams: predicted probability on the x-axis, empirical disrupted-rate on the y-axis. Closer to the diagonal is better. Scenario-B XGBoost and LightGBM are tightly calibrated; logistic regression is dramatically over-confident at the top of the probability range.}
    \label{fig:calibration}
\end{figure}

\subsection{Per-country, per-position, per-train-class breakdowns}

Aggregate test-set PR-AUC hides substantial structural variation. Scenario-B XGBoost reaches PR-AUC \num{0.989} on Finland (\num{495000} test rows), \num{0.858} on the Netherlands (\num{1790000} rows) and only \num{0.788} on Italy (\num{21000} rows). \Cref{fig:per-country} shows the per-country breakdown; \cref{fig:per-position} faceted by stop position quartile; \cref{fig:per-class} faceted by train class.

\begin{figure}[htbp]
    \centering
    \includegraphics[width=0.92\textwidth]{stop_level/per_country_metrics.png}
    \caption{Test-set PR-AUC by country, faceted by scenario. Finland's small station count and high disruption rate make the per-station lag features unusually informative there; Italy is the hardest target, partly due to its small test sample and partly due to a known weather-feature bug discussed in \cref{sec:transfer}.}
    \label{fig:per-country}
\end{figure}

\begin{figure}[htbp]
    \centering
    \includegraphics[width=0.92\textwidth]{stop_level/per_position_metrics.png}
    \caption{PR-AUC by stop position quartile. In scenario~A the model is weakest near the origin (where less history exists). In scenario~B the curve flattens: the prior-stop signal dominates regardless of route position.}
    \label{fig:per-position}
\end{figure}

\begin{figure}[htbp]
    \centering
    \includegraphics[width=0.92\textwidth]{stop_level/per_train_class_metrics.png}
    \caption{PR-AUC by train class. Class~0 is local / commuter, class~5 is ultra-high-speed. The high-speed and intercity classes (3, 4, 5) are easier than local services, which have more variable headways and fewer dedicated tracks.}
    \label{fig:per-class}
\end{figure}

\subsection{Cascading respect}

A useful sanity check on scenario~B is to ask whether the model actually \emph{uses} the inflight signal it is given. For each model we compute the predicted-disruption probability conditional on the prior stop being late versus on time. A meaningful gap means the model is reading the cascade. \Cref{fig:cascading} shows the result: the three tree-based models give roughly \num{0.97} when the prior stop is late versus \num{0.05} when it is on time, while logistic regression gives a softer \num{0.65} versus \num{0.05}.

\begin{figure}[htbp]
    \centering
    \includegraphics[width=0.78\textwidth]{stop_level/cascading_respect.png}
    \caption{Scenario~B only. For each model we plot $P(\text{disrupted}\mid\text{prior stop late})$ against $P(\text{disrupted}\mid\text{prior stop on time})$. A wide gap (top-left) confirms the model is reading the cascading delay; the tree-based models all lie there.}
    \label{fig:cascading}
\end{figure}
